\section*{Chapter \ref{ch:mx:portfolio-trades-and-risk-bids} --- Portfolio Trades and Risk Bids}

\begin{solution}{exo:mx:portfolio-trades-and-risk-bids:1}
$0.3/0.2=1.5$ days; $3+0.7\times200\times\sqrt{0.2}=65.6$ basis points.
\end{solution}

\begin{solution}{exo:mx:portfolio-trades-and-risk-bids:2}
The variance is $\sigma^2d/3$: $150\sqrt{3/3}=150$ basis points.
\end{solution}

\begin{solution}{exo:mx:portfolio-trades-and-risk-bids:3}
$4\times\E[\max\text{ of 6 normals}]=4\times1.27=5.07$ basis points.
\end{solution}

\begin{solution}{exo:mx:portfolio-trades-and-risk-bids:4}
A long basket of 500 stocks is mostly market risk, which one futures trade removes at once. What remains is sector and style risk (the other
factors) and the stocks' specific risk, which the future does not hedge; with 500 names the specific part is small, so the other factors dominate.
\end{solution}

\begin{solution}{exo:mx:portfolio-trades-and-risk-bids:5}
The fund held more of the liquid names within each bucket than a random draw from the bucket does; the profile cannot show it. A basket tilted to the
less liquid names of each bucket would be underestimated, which is the blind bidder's real danger.
\end{solution}

\begin{solution}{exo:mx:portfolio-trades-and-risk-bids:6}
The errors become correlated: the common part of the error is not selected by the auction (every bidder shares it), so the curse from the independent
part shrinks, but all bidders are wrong together, and the client gains or loses from the shared error.
\end{solution}

\begin{solution}{exo:mx:portfolio-trades-and-risk-bids:7}
Cost 14.9, risk charge 25.0, shading 6.2 (the estimation error scales with the larger base), margin 2.0: 48.1 basis points.
\end{solution}

\begin{solution}{exo:mx:portfolio-trades-and-risk-bids:8}
Winning says the bids were low relative to the competitors, not that they covered the costs: the winner's curse selects the trades whose costs were
underestimated. Only the realised P\&L of the won trades against their bids says whether the pricing was right.
\end{solution}

\begin{solution}{pb:mx:portfolio-trades-and-risk-bids:1}
\textbf{1.} Agency: the client bears costs and risk; principal: the broker takes the positions at an agreed price.
\textbf{2.} See the definitions.
\textbf{3.} 786 names, median daily volume USD~37.5 million, median volatility 171 basis points; a USD~250 million basket of 500 names, median
position 1.0\% of daily volume.
\textbf{4.} A transaction-cost saving for the client and an efficient source of flow for liquidity providers.
\textbf{5.} At most 20\% of daily volume, half-spread plus $Y\sigma\sqrt{\min(Q/V,0.2)}$; 14.9 basis points.
\textbf{6.} $\int_0^a(1-t/a)(1-t/b)\,dt=a-a/2-a^2/(2b)+a^2/(3b)=a/2-a^2/(6b)$.
\textbf{7.} 40.4 basis points unhedged, 12.5 hedged.
\textbf{8.} 98.8\% of the notional is below 5\% of daily volume; no position needs more than a day.
\textbf{9.} The lowest of four independent estimates is too low by $1.03s$ on average; shade by that.
\textbf{10.} $-2.2$ basis points without shading, 2.0 (the margin) with it.
\textbf{11.} By pricing 300 baskets with the same profile: cost 17.0 (s.d. 0.4), hedged risk 13.0 (1.1).
\textbf{12.} Disclosed: 14.9, 12.5, 4.2 and 2.0; blind: 17.0, 13.0, 4.4 and 2.0.
\textbf{13.} \emph{Named result}: 33.6 basis points below the close for the disclosed basket and 36.4 for the blind one; the winner's-curse shading
against three competitors is 4.2 and 4.4 basis points.
\textbf{14.} 50.4, 32.8 and 24.0 basis points of the fund.
\textbf{15.} It pays no spread and moves no price: both sides are satisfied at the mid.
\textbf{16.} When the names would leak to dealers who would trade ahead of the unwind by more than the premium.
\textbf{17.} The impact model's coefficient and the competitors' errors being independent.
\textbf{18.} Shade less for the shared error, but widen the margin for model risk.
\textbf{19.} Principal buys certainty at a price; agency is cheaper on average when the client can bear the risk.
\textbf{20.} The certainty of the close, for the cost and risk of the unwind and the winner's curse.
\end{solution}

\begin{solution}{iq:mx:portfolio-trades-and-risk-bids:1}
The names (or profile), sizes against volume, spreads and volatilities, the covariance and available hedges, the benchmark, the competition. Price
= expected liquidation cost + a charge for the unwind's risk + winner's-curse shading + margin.
\iqlookfor{The four parts; hedging.}
\end{solution}

\begin{solution}{iq:mx:portfolio-trades-and-risk-bids:2}
The winning estimate is the most optimistic of several noisy estimates; correct by adding the expected lead of the best of $n$ estimates, $s\E[\max]$,
or by pricing conditional on winning.
\iqlookfor{Selection on the lowest estimate.}
\end{solution}

\begin{solution}{iq:mx:portfolio-trades-and-risk-bids:3}
The variance of the P\&L over the unwind, $\sum w_iw_jC_{ij}\int x_ix_j\,dt$ with a risk model; hedge the market (and sectors or styles) with futures
or ETFs at once, leaving specific and unhedgeable factor risk.
\iqlookfor{Time-weighted covariance; hedges.}
\end{solution}

\begin{solution}{iq:mx:portfolio-trades-and-risk-bids:4}
Whether the profile hides adverse tilts (the client's baskets are less liquid within buckets than assumed), whether the shading reflects the real
number of competitors and the model's error, and whether the losses come from cost, risk or leakage after winning.
\iqlookfor{Profile bias; shading; decomposition of the P\&L.}
\end{solution}

\begin{solution}{iq:mx:portfolio-trades-and-risk-bids:5}
Inventory both portfolios, net the common names, cross internally, hedge the out-of-market exposure with futures, schedule the rest by liquidity, and
report the cost against the pre-trade estimate.
\iqlookfor{Netting, crossing, hedging, reporting.}
\end{solution}

\begin{solution}{iq:mx:portfolio-trades-and-risk-bids:6}
Draw baskets consistent with the profile and price them; check the bias by pricing, after the fact, disclosed baskets as if they were blind and
comparing with their true cost over many trades.
\iqlookfor{Simulation from the profile; back-testing the blind estimate.}
\end{solution}
